\documentclass[12pt,a4paper]{article}

\usepackage[utf8]{inputenc}
\usepackage[T1]{fontenc}
\usepackage[margin=2.5cm]{geometry}
\usepackage{hyperref}
\usepackage{microtype}
\usepackage{parskip}
\usepackage{lmodern}

\hypersetup{colorlinks=true, urlcolor=blue!60!black}

\pagestyle{empty}

\begin{document}

\begin{flushright}
Kundai Farai Sachikonye \\
Auf der Lichtung 19, 82178 Puchheim, Germany \\
\href{https://github.com/fullscreen-triangle}{github.com/fullscreen-triangle} \\
\texttt{kundai.sachikonye@bitspark.com} \\
(+49) 1626386344 \\[6pt]
\today
\end{flushright}

\vspace{1em}

Fraunhofer-Institut für Intelligente Analyse- und Informationssysteme IAIS \\
Abteilung NetMedia \\
Dresden

\vspace{1.5em}

\textbf{Re: Research Engineer GenAI --- Post-Training and Agentic AI}

\vspace{1em}

Dear Hiring Team,

I am applying for the Research Engineer position in GenAI with a focus on
post-training and agentic AI. Two parts of the role line up unusually well with
work I have already built: agentic / multi-agent systems, and foundation-model
adaptation --- and the role's emphasis on abstract, conceptual work is exactly how
I prefer to work.

\textbf{Agentic and multi-agent AI.}
Rather than only using agent frameworks, I have built one. Kwasa-kwasa is a
declarative orchestration language for multi-agent, evidence-bearing computation:
a runtime that delegates sub-computations to multiple substrates (LLM oracles,
external tools, Python specialists), aggregates graded confidence, and grounds
generated values against observation before they are committed. It has a Rust core
and a TypeScript runtime and ships with VSCode tooling including a language server
(\href{https://github.com/fullscreen-triangle/kwasa-kwasa}{github.com/fullscreen-triangle/kwasa-kwasa}).
Alongside it, four-sided-triangle is an eight-stage metacognitive RAG / reasoning
pipeline with resource-aware orchestration. Designing multi-agent systems and
moving them toward production is the work I most want to do.

\textbf{Foundation-model post-training.}
On the model side, I work in PyTorch and have done foundation-model adaptation
through low-rank (LoRA) fine-tuning and knowledge distillation, with a research
interest in training methodology --- specifically curriculum design and inverse
curricula. I am candid that the specific RL-based post-training methods and
frameworks you name (DPO, GRPO, continual pre-training, NeMo-RL, Megatron) are
ones I would ramp up on rather than ones I have run at cluster scale; the
foundations --- deep ML understanding, PyTorch, SLURM-based GPU training, and
building evaluation/benchmarking pipelines --- are already in place.

\textbf{Research that ships.}
Everything I build I take to a deployed, documented, benchmarked state, validated
against reference data. Turning research methods into scalable, production-near
systems is the default of how I work, which is the core of this role.

I would welcome the chance to discuss how my agentic-AI and foundation-model work
could contribute to NetMedia's projects.

\vspace{1.5em}
Yours sincerely,

\vspace{2.5em}
\textbf{Kundai Farai Sachikonye}

\end{document}
